
\documentclass[10pt, openany]{book}

\usepackage{../preamble}
\addbibresource{c:/Dev/CQT/ref.bib}

\fancyhead[LE,RO]{\thepage}
\fancyhead[RE]{\textbf{Collection}}
\fancyhead[LO]{\textit{Math \& Physics}}

\begin{document}


\frontmatter
\title{\textbf{Collection}\\ \large Math \& Physics}
\author{Lao Chen}
\date{\today}
\maketitle


\begin{center}
\large
专栏链接：\url{https://www.zhihu.com/column/ps-math}\\[6pt]
作者主页：\url{https://www.zhihu.com/people/smdw}\\[6pt]
文章总数：约180篇，编号体系 MP1 $\sim$ MP149+\\[6pt]
时间跨度：2018 $\sim$ 2025
\end{center}

\tableofcontents

\mainmatter

\chapter{微分流形基础}

\section{MP10：微分流形上的函数和参数曲线}
\label{sec:mp10}

\href{https://zhuanlan.zhihu.com/p/44203444}{原文链接}

本讲是为后面在微分流形上引入余切空间和切空间做技术准备。主要的技术是方向导数。由于微分流形不同于欧氏空间，我们无法直接定义导数，需要借助函数和参数曲线的复合映射来实现微分运算。

在微分流形 $M$ 上，光滑函数 $f: M \to \mathbb{R}$ 构成一个代数结构。参数曲线 $\gamma: \mathbb{R} \to M$ 则提供了对流形进行「探测」的手段。通过函数与参数曲线的复合 $f \circ \gamma: \mathbb{R} \to \mathbb{R}$，我们将问题拉回到熟悉的实数域，从而可以在流形上定义方向导数。

\medskip
\noindent\rule{\textwidth}{0.4pt}

\section{MP11：余切空间}
\label{sec:mp11}

\href{https://zhuanlan.zhihu.com/p/44190499}{原文链接}

本讲介绍微分流形一点上的余切空间。

\textbf{函数芽}：局部化等价。对函数芽求方向导数，引出余切空间。利用参数曲线构造复合映射，对参数求复合导数，类似方向导数的运算。

我们的注意力只在函数在无限接近一点时的性质，是一种局部化的处理。在谈论一点附近的函数性质时，用函数芽代替函数。函数芽的集合上可以定义加法和数乘运算，构造线性空间的结构。这里的加法和数乘相当于函数上相同运算后的等价化（局部化）。

进一步考察方向导数构造的某种映射的核，可以验证它构成一个线性子空间，从而可以构造商空间。

\medskip
\noindent\rule{\textwidth}{0.4pt}

\section{MP12：切空间}
\label{sec:mp12}

\href{https://zhuanlan.zhihu.com/p/44211005}{原文链接}

这一讲继续余切空间的讨论，介绍切空间。

切空间是余切空间的对偶空间，切向量是余切向量的对偶向量。在微分流形 $M$ 的一点 $p$ 处，余切空间 $T_p^*M$ 由函数芽的微分构成，而切空间 $T_pM$ 则定义为余切空间的对偶：
\[
T_pM = (T_p^*M)^*
\]
切向量可以自然地理解为作用在函数上的方向导数算子。

\medskip
\noindent\rule{\textwidth}{0.4pt}

\section{MP19：拓扑加油站(1)：道路连通、同伦、基本群}
\label{sec:mp19}

\textit{（编号和标题来自专栏导读，原文链接待确认）}

从连通性开始，不连通的拓扑空间可以通过连通分支分解，每个分支的开集组合仍保持开集性质。道路连通是比连通更强的条件。

\textbf{同伦}：两个连续映射 $f, g: X \to Y$ 之间的同伦是一个连续映射 $H: X \times [0,1] \to Y$，满足 $H(x,0) = f(x)$，$H(x,1) = g(x)$。道路同伦则要求基点保持不动。同伦是等价关系。

\textbf{基本群} $\pi_1(X, x_0)$：以 $x_0$ 为基点的道路同伦等价类构成的群。基本群是一个拓扑不变量，用于描述拓扑空间的「洞」的数量。在同伦意义下，基本群在同伦等价下不变。

\medskip
\noindent\rule{\textwidth}{0.4pt}

% ============================================================


% 参考文献
%\bibliographystyle{apalike}
%\bibliography{ref}
\printbibliography

\end{document}